\chapter{Architecture: the zones and their boundaries}
\label{ch:architecture}

\autoref{ch:philosophy} ended on a principle and an unpaid debt. The principle was
\emph{dissolve, do not guard}: sanctity at a boundary is achieved by removing the
need for the dangerous capability, not by permitting it under controls
(\autoref{sec:sacred}). The debt is the obvious follow-up question, and it is an
engineering question, not a philosophical one --- \emph{where do you put the
boundary so that the dangerous permission stops being necessary?}

This chapter is the answer, and it has a shape worth stating before the details
arrive. The same move is made five times, at five different scales: between
processes, between operating-system users, between modules, between types, and
between the terms of a single query. Each time the move is identical --- take a
capability that a rule was being asked to restrain, and relocate the boundary until
the capability is \emph{absent} on one side of it. What differs is only the
enforcement machinery available at that scale, and the honesty with which the
design admits how strong its machinery actually is.

That reading is the manual's own synthesis and is flagged as such; but it is not
invented out of nothing. The record states the correspondence explicitly for two of
the five scales: the foundation denylist is ``ungrantable at the centre with
grantable strata around it,'' and the code side ``now has the same geometry, with
\texttt{CONSTITUTION.md} to data what the inner ring is to machinery: the fixed
points'' (\artifact{dn-inner-outer-core}, \S 2.8). What follows extends that
observation across the remaining three.

\section{The grading rubric: three tiers of enforcement}
\label{sec:tiers}

Before drawing a single boundary it is worth fixing how the chapter will judge one,
because a design manual that describes only mechanisms is a catalogue, and a
catalogue cannot tell you whether a system is safe.

The record supplies the rubric. Enforcement sits on a three-rung ladder ---
$\tstruct{} \succ \tstatic{} \succ \tconv{}$ --- and the ladder carries one
non-obvious composition law: \emph{the tier of a composed grant is the minimum tier
along the construction chain} (\artifact{dn-capability-scope}, \S 2.2). The law is
what stops a convention-tier permission from laundering itself through a
structural-looking wrapper and emerging apparently strong. Its practical
consequence is that the tier is an \emph{annotation on the whole chain}, never a
property of the last link, and it must be recorded honestly at each link or the
minimum cannot be computed at all.

\begin{principle}[Enforcement has a tier, and the tier is the minimum]
Every boundary in this system is labelled with the strength of the mechanism that
holds it: \tstruct{} (violation is unrepresentable or refused by the kernel),
\tstatic{} (violation is caught mechanically before the program runs), or
\tconv{} (violation is caught by review, or not at all). A boundary composed of
several mechanisms inherits the weakest of them
(\artifact{dn-capability-scope}, \S 2.2).
\end{principle}

\autoref{fig:enforcement-ladder} shows the ladder with the system's single most
enforced invariant --- the sealed core's network isolation --- placed on all three
rungs at once, which is exactly the point: the tiers stack rather than replace.
\autoref{tab:five-scales} is then the chapter's map, and it doubles as an honest
scorecard. Two rows are \tconv{}, and the chapter says so where it reaches them.

\begin{figure}[t]
  \centering
  % figures/enforcement-ladder.tex — one invariant, three tiers of enforcement.
% The LADDER (structural > static+guard > convention, and "composition takes the min
% tier along the construction chain") is \artifact{dn-capability-scope} §2.2.
% The three mechanisms shown are the three enforcements of the SAME property
% (non-negotiables 1 and 2, \coderef{docs/BUILD-SPEC.md}{\gitref} §3):
%   runtime  — \coderef{core/sealing.py}{\gitref} (fail-closed in-process guard)
%   static   — \coderef{ops/import\_lint.py}{\gitref} (the import-closure lint)
%   kernel   — \coderef{ops/network/ouroboros-egress.pf.conf}{\gitref} (pf on uid),
%              per \artifact{dn-plane-principals} §3.4.
% "The import firewall does not retire — it stays as the cheap, pre-runtime tripwire
% (defense in depth)" is \artifact{dn-plane-principals} §3.4, verbatim in substance.
% \input by chapters/02-architecture.tex inside a figure environment.
\begin{tikzpicture}[
    rung/.style={draw=inkblue, thick, rounded corners=2pt, fill=panel,
      align=center, inner sep=4pt, text width=30mm, minimum height=12mm,
      font=\small\sffamily},
    mech/.style={align=left, font=\scriptsize\sffamily, text width=62mm},
    when/.style={align=center, font=\scriptsize\itshape, text=soft, text width=22mm},
    up/.style={-{Stealth[length=2mm]}, very thick, gate},
  ]
  \node[rung] (conv) at (0,0) {\tconv{}};
  \node[rung] (stat) at (0,17mm) {\tstatic{}};
  \node[rung] (strc) at (0,34mm) {\tstruct{}};

  \draw[up] (conv.north) -- (stat.south);
  \draw[up] (stat.north) -- (strc.south);
  \node[font=\scriptsize\itshape, text=gate, right=1mm, rotate=90,
        anchor=south] at ($(conv.north)!0.5!(stat.south)$) {stronger};

  \node[when, left=4mm of conv] {caught at review, or not at all};
  \node[when, left=4mm of stat] {caught before the program runs};
  \node[when, left=4mm of strc] {cannot happen};

  \node[mech, right=5mm of conv]
    {``core must not open a socket'' --- written in a docstring and believed.};
  \node[mech, right=5mm of stat]
    {\texttt{ops/import\_lint.py} --- no module under \texttt{core/} \emph{reaches} a
     network-capable module, proved by reading the AST.};
  \node[mech, right=5mm of strc]
    {\texttt{core/sealing.py} --- a fail-closed process-wide guard (runtime);\\
     the \texttt{pf} anchor on the \pcore{} uid --- the kernel refuses the packet.};

  \node[font=\scriptsize\itshape, text=inkblue, align=left, text width=100mm,
        anchor=north west] at ([shift={(-30mm,-8mm)}]conv.south west)
    {The tiers stack rather than replace: the static lint is retained as the cheap
     pre-runtime tripwire even after the kernel fact exists. And a grant composed
     through several mechanisms inherits the \emph{minimum} tier along the chain ---
     which is what stops a convention-tier permission laundering through a structural
     wrapper.};
\end{tikzpicture}

  \caption[The enforcement ladder]{One invariant, three tiers. The rungs are the
    ladder of \artifact{dn-capability-scope} \S 2.2; the mechanisms are the three
    enforcements the repository actually carries for non-negotiables 1 and 2
    (\coderef{docs/BUILD-SPEC.md}{\gitref}, \S 3). The static lint is deliberately
    \emph{retained} after the kernel fact exists --- ``it stays as the cheap,
    pre-runtime tripwire; it catches the violation at review, the kernel catches it
    at execution'' (\artifact{dn-plane-principals}, \S 3.4).}
  \label{fig:enforcement-ladder}
\end{figure}

\begin{table}[t]
  \centering\footnotesize
  \begin{tabular}{@{}llll@{}}
    \toprule
    \textbf{Scale} & \textbf{The boundary} & \textbf{Held by} & \textbf{Tier} \\
    \midrule
    processes  & zones A/B/C           & a filesystem handoff        & \tstruct{} \\
    modules    & the import closure    & an AST lint                 & \tstatic{} \\
    OS users   & four plane principals & kernel DAC + \texttt{pf} & \tstruct{} \\
    modules    & inner ring / outer    & a computed fixed point      & \tstatic{} \\
    types      & the checked region    & a strict type checker       & \tstatic{} \\
    queries    & the scope lattice     & construction-time admissibility & mixed \\
    workflow   & the plan's write scope & pre-hoc guard + post-hoc audit & mixed \\
    \bottomrule
  \end{tabular}
  \caption[The five scales]{The chapter's map. ``Mixed'' is not evasion: the scope
    lattice labels each of its five inhabitants individually
    (\artifact{dn-capability-scope}, \S 2.4), and the workflow's write enforcement
    is deliberately porous pre-hoc and tight post-hoc
    (\artifact{dn-agent-workflow}, \S 6). \autoref{sec:residuals} states every
    residual.}
  \label{tab:five-scales}
\end{table}

\section{Zones: the boundary is a spool of files}
\label{sec:zones}

The outermost boundary is drawn between \emph{processes}, and it is drawn around
the network.

Three zones (\coderef{docs/BUILD-SPEC.md}{\gitref}, \S 6). \zoneA{}, the sealed
core, holds everything private: the vault, the vector store over the thought-graph,
the telemetry store, the private-reasoning models, the librarian, curator, dreaming
and matching agents, the agent factory, the sandbox broker. It has no outbound
network. \zoneB{}, the networked edge, is two narrowly scoped processes --- a
bridge that ships research specs outward and retrieves public corpora, and an
interface gateway that relays the owner's messages --- each running with the
private vault \emph{unmounted}. \zoneC{} is the cloud: encrypted backups and
outbound research fetchers, which see only de-identified topic criteria and public
literature.

\begin{figure}[t]
  \centering
  % figures/trust-zones.tex — the three trust zones and the airlock asymmetry.
% Source: \coderef{docs/BUILD-SPEC.md}{\gitref} §6 (zones A/B/C; the boundary is a
% filesystem handoff, not shared memory or shared DB access; the data-flow asymmetry)
% and §3 invariants 1, 2, 11. The handoff-as-airlock reading is
% \artifact{dn-plane-principals} §3.1 ("the only thing edge and core share is a spool
% of typed files"). Zone names from notation.tex.
% \input by chapters/02-architecture.tex inside a figure environment.
\begin{tikzpicture}[
    node distance=8mm and 10mm,
    zone/.style={draw=inkblue, thick, rounded corners=3pt, fill=panel,
      align=center, inner sep=4pt, text width=31mm, minimum height=20mm,
      font=\small\sffamily},
    spool/.style={draw=gate, thick, fill=white, align=center, inner sep=3pt,
      text width=15mm, minimum height=13mm, font=\scriptsize\sffamily},
    outb/.style={-{Stealth[length=2mm]}, thick, gate},
    inb/.style={-{Stealth[length=2mm]}, thick, inkblue},
    lbl/.style={font=\scriptsize\itshape, align=center, text=soft},
  ]
  \node[zone] (a) {\zoneA{} --- sealed core\\[2pt]
    \scriptsize vault · vector store · private models · the agents\\[2pt]
    \scriptsize\bfseries no outbound network};
  \node[spool, right=of a] (h) {handoff\\directories\\[2pt]\itshape typed files};
  \node[zone, right=of h] (b) {\zoneB{} --- networked edge\\[2pt]
    \scriptsize bridge · interface gateway\\[2pt]
    \scriptsize\bfseries vault unmounted};
  \node[zone, right=of b] (c) {\zoneC{} --- cloud\\[2pt]
    \scriptsize encrypted backups ·\\ outbound research fetchers};

  \draw[outb] (a.east) -- node[lbl, above] {sanitized\\job specs} (h.west);
  \draw[inb]  (h.east) -- node[lbl, above] {results} (b.west);
  \draw[outb] ([yshift=3mm]b.east) -- node[lbl, above] {de-identified\\topic criteria} ([yshift=3mm]c.west);
  \draw[inb]  ([yshift=-5mm]c.west) -- node[lbl, below] {public\\content only} ([yshift=-5mm]b.east);

  % The one boundary that is a *fact about processes*, not a convention.
  \draw[gate, thick, dashed] ($(h.north)+(0,7mm)$) -- ($(h.south)+(0,-9mm)$);
  \node[lbl, text=gate, below=9mm of h.south, text width=30mm]
    {the boundary is a \textbf{filesystem handoff} --- never shared memory,\\ never a shared DB};
\end{tikzpicture}

  \caption[The trust zones and the airlock]{The three zones, the handoff between
    them, and the airlock's deliberate asymmetry --- de-identified criteria out,
    public content in. Source: \coderef{docs/BUILD-SPEC.md}{\gitref}, \S 6 and
    \S 3 (invariants 1, 2 and 11); the reading of the handoff directory as an
    airlock is \artifact{dn-plane-principals}, \S 3.1.}
  \label{fig:trust-zones}
\end{figure}

The load-bearing sentence in that section is not the zone list --- it is the one
about what connects them. \emph{The boundary between the core and the networked
processes is the filesystem handoff (specific directories), not shared memory or
shared DB access} (\coderef{docs/BUILD-SPEC.md}{\gitref}, \S 6). This looks like an
implementation detail and is in fact the whole design. A shared database or a
shared memory segment is a channel with \emph{ambient} authority: whatever is
reachable through it is reachable by anyone holding the handle, and the only thing
standing between a compromised edge process and the corpus is the correctness of
every query the edge is trusted not to write. A directory of typed files has no
ambient authority at all. It carries exactly what was written into it, it is
inspectable after the fact, and --- the property that matters most --- it is
\emph{inert}: reading a file is not an invocation. When the plane split later
formalises this, it names the handoff for what it always was: ``the only thing edge
and core share is a spool of typed files --- no shared network, no shared uid''
(\artifact{dn-plane-principals}, \S 3.1).

The second load-bearing idea is the airlock's asymmetry. Outbound traffic is
de-identified topic criteria only; inbound traffic is public content only
(\coderef{docs/BUILD-SPEC.md}{\gitref}, \S 6). The point is not symmetry-breaking
for its own sake. It is that these two constraints, taken together, let the system
keep learning from the world without the world ever learning about the owner: ``the
owner's inner life never leaves the walls; fresh external knowledge still gets
in.'' The same section adds the warning that shows the design thought past its own
slogan --- \emph{even topic keywords carry intent}, so outbound criteria are kept
general. A boundary that leaks the shape of a question has leaked something.

One further asymmetry is stated at the constitutional level and belongs here
because it is the only place the design knowingly lets something cross: the
interface may transit a third party; the corpus never does
(\coderef{docs/BUILD-SPEC.md}{\gitref}, \S 3, invariant 11). If a messaging adapter
routes through an external service, the owner's \emph{interactions} leave the
boundary even though the corpus does not --- and the record insists that this
tradeoff be explicit and opt-in, with the private default keeping interactions
inside the walls. It is worth pausing on the register of that sentence. It does not
claim the boundary is perfect; it names the one place it is deliberately porous and
requires the owner to choose it.

\section{The import firewall: turning a zone into a theorem}
\label{sec:firewall}

A zone drawn on an architecture diagram is a convention --- the lowest rung of
\autoref{fig:enforcement-ladder}. Nothing about \zoneA{} being labelled ``no
network'' prevents a module inside it from importing \texttt{urllib} on a Tuesday.
The question this section answers is how the label becomes a fact, and the answer
contains the single most elegant argument in the codebase.

\subsection*{The idea: make it a property of the closure, not of a module}

Checking that no core module opens a socket is a whack-a-mole exercise: you check
every module, and you must re-check on every change, and a module that merely
\emph{imports} something that opens a socket passes the check while remaining
capable. The move is to stop reasoning about modules and start reasoning about the
\emph{import closure}. The lint's own docstring states the argument
(\coderef{ops/import\_lint.py}{\gitref}): if no module under \texttt{core/} can
\emph{reach} a network-capable module, then no egress path exists regardless of
edge behaviour --- \emph{composition cannot create one}.

That is worth making precise, because the precision is where the strength comes
from --- and because the precision also exposes exactly what the lint does
\emph{not} yet buy. Let $G$ be the module import digraph of the repository, with an
edge $m \to m'$ whenever $m$ imports $m'$; write $\imports(m)$ for the direct
imports of $m$ and $\reach(S)$ for everything reachable from $S$ by any number of
edges. Let $\coremods{}$ be the modules under \texttt{core/}, let $\netmods{}$ be
the network-capable module roots, and let $\netallow{} \subseteq \coremods{}$ be the
audited allowlist.

What the lint actually checks is \emph{local} --- one AST walk per file, direct
imports only (\coderef{ops/import\_lint.py}{\gitref}):

\begin{invariant}[Zone non-interference and network containment]
\label{inv:firewall}
For every $m \in \coremods{}$: $\imports(m) \cap \{\texttt{edge},\texttt{cloud}\}
= \varnothing$, and if $m \notin \netallow{}$ then
$\imports(m) \cap \netmods{} = \varnothing$
(\coderef{ops/import\_lint.py}{\gitref}; \coderef{docs/BUILD-SPEC.md}{\gitref},
\S 3, invariants 1--2).
\end{invariant}

\noindent A local check cannot, by itself, give a global reachability property. What
it gives is this:

\begin{proposition}[Every escape route is visible]
\label{prop:closure}
Under \autoref{inv:firewall}, every path in $G$ from a module in $\coremods{}$ to a
module in $\netmods{}$ either begins at a module of $\netallow{}$, or leaves
$\coremods{}$ at some step into a module that is neither \texttt{edge},
\texttt{cloud}, nor network-capable.
\end{proposition}

\begin{proof}
Let $m_0 \to \cdots \to m_k$ be such a path with $m_0 \in \coremods{}$ and
$m_k \in \netmods{}$, and let $i$ be the last index with $m_i \in \coremods{}$. If
$i = k$ then $m_k \in \coremods{} \cap \netmods{}$, which \autoref{inv:firewall}
permits only for $m_k \in \netallow{}$. Otherwise $m_{i+1} \in \imports(m_i)$ lies
outside $\coremods{}$, and \autoref{inv:firewall} forbids it from being
\texttt{edge}, \texttt{cloud}, or a member of $\netmods{}$ unless
$m_i \in \netallow{}$.
\end{proof}

\noindent So the local condition closes every \emph{direct} route and forces every
remaining route to pass through a module core imports from outside itself. The
global claim --- \emph{no egress path exists, regardless of composition} --- therefore
holds under one extra hypothesis:

\begin{quote}
$(\star)$\quad every module core reaches outside $\coremods{}$ is stdlib or a pinned,
side-effect-free third-party library.
\end{quote}

And that hypothesis is not decoration; it is a second, separately named invariant
with its own enforcement, and \textbf{it is currently false on purpose}. Core still
imports from sibling packages, the audit that measured it is
\artifact{finding-0103}, and the test asserting $(\star)$ is \emph{red by design} ---
a ratchet whose count may only ever decrease, which the owner chose over a silent
allowlist or a deferred cleanup
(\coderef{tests/unit/test\_core\_self\_containment.py}{\gitref};
\artifact{dn-inner-outer-core}, \S 2.4 records it counting 19 down to 0).

\begin{proposition}[The closure claim, at the limit]
\label{prop:closure-limit}
If $(\star)$ holds, then
$\reach(\coremods{}) \cap \netmods{} \subseteq \netallow{}$: composition cannot
create an egress path.
\end{proposition}

\begin{proof}
By \autoref{prop:closure}, a path from core to $\netmods{}$ that does not begin in
$\netallow{}$ must leave core into some $m_{i+1} \notin \coremods{}$. Under
$(\star)$ that $m_{i+1}$ is stdlib or a pinned pure library, none of which is in
$\netmods{}$ or reaches it --- so no such path exists.
\end{proof}

This is the chapter's first real piece of synthesis, and it is worth stating
plainly: \emph{the import firewall and the core self-containment ratchet are not two
hygiene programmes; they are the two hypotheses of one theorem}. Until the ratchet
reaches zero, the sealed core's static isolation is a conditional result, and the
condition is measurable --- which is a far more useful state to be in than either
overclaiming the guarantee or not knowing what it depends on. It is also why the
runtime guard and the kernel anchor of the next section are not redundancy: they
hold the property unconditionally while the static argument is still being
discharged.

The lint's own contribution remains large and cheap. It needs no understanding of
the program --- it reads imports --- and the record names the move: a runtime
guarantee lifted to a static one, ``provable by reading the AST without running the
program'' (\artifact{dn-type-system-as-core-audit}, \S 2.4).

\subsection*{Two rules of different strength --- and the honesty about the second}

The firewall is not one rule but two, and they are deliberately unequal.

\begin{quote}\small\ttfamily
source: ops/import\_lint.py@\gitref
\end{quote}
{\small\begin{verbatim}
FORBIDDEN_ZONES: frozenset[str] = frozenset({"edge", "cloud"})

    ... (NETWORK_MODULES, the audited network-capable set, elided) ...

NETWORK_ALLOWLIST: frozenset[str] = frozenset({
    "core/sealing.py",                 # wraps socket to seal it
    "core/models/ollama_client.py",    # loopback Ollama IPC channel (127.0.0.1)
})
\end{verbatim}}

\noindent The first rule --- no \texttt{core/} module may import \texttt{edge} or
\texttt{cloud} --- is hard, with zero exceptions, and the docstring explains why it
can afford to be: the sealed core ``cannot even \emph{name} the networked zones, so
there is no \texttt{core} $\to$ \texttt{edge} $\to$ net path to create.'' This is
capability dissolution at the module scale. There is nothing to guard because there
is nothing to reference.

The second rule --- no networking primitive --- carries an allowlist of exactly two
audited files, and the design refuses to overstate what that buys. The docstring's
own words: the guarantee is ``exactly these two audited files touch networking
primitives, both loopback-only,'' \emph{not} ``no core file imports \texttt{http}''
--- and overclaiming the latter ``would be dishonest (the Ollama channel is
real).'' A reader should notice what has happened there. The system needs to speak
to a locally resident model server, so a loopback channel genuinely exists; rather
than hide it behind a stronger-sounding claim, the design names it, audits it,
pins it to two files, and writes the weaker true sentence into the source. That
habit --- preferring the weaker claim that is true --- is the same stance
\autoref{sec:falsifiers} identified in the process, now visible in a docstring.

\subsection*{The same invariant, three times}

The import lint is the middle rung. Beneath it, at runtime,
\coderef{core/sealing.py}{\gitref} installs a process-wide, fail-closed guard that
intercepts every socket \texttt{connect} and permits only loopback destinations and
\texttt{AF\_UNIX}, raising before a packet leaves. Its docstring is candid about
its own ceiling --- a native extension opening its own socket could bypass a
Python-level guard --- and about why it exists anyway: it fails closed, it is
process-wide, it is verifiable in tests, and it composes with the OS-level layer.
Above the lint, at the kernel, sits the \texttt{pf} anchor of
\autoref{sec:planes} --- committed as an inert file, loaded by the owner as part
of a deliberate migration. Three mechanisms, one property, three tiers, each with
its stated limit. That is what defence in depth looks like when the layers are
honest about their individual weakness.

\section{Planes: the boundary the kernel enforces}
\label{sec:planes}

The zone boundary is about \emph{code}. The next boundary is about \emph{processes
and the users that own them} --- and it is the one place where the system's rules
stop being the system's rules and become the operating system's.

\begin{devolution}
The design arrived here in two steps, and the first is retained because the second
is unintelligible without it. \artifact{dn-ouroboros-principal}
(\textbf{superseded}) minted \emph{one} new principal --- \texttt{ouroboros}, the
live system --- and drew the boundary between ``the human builds the machine'' and
``the machine is running.'' Its \S 3.3 stated the generalisation it did not itself
take: \emph{the uid separates planes, not agent types}. Then \artifact{finding-0116}
observed that exhaust has two producer planes --- the workflow writes build
reports, the system writes its own emissions --- so one new principal cannot carry
the plane model the note had articulated. \artifact{dn-plane-principals}
supersedes it on that warrant, generalising rather than patching: \emph{four OS
principals, one per plane}. The superseded note's threat model, role-account
mechanics, daemon relocation and execution protocol are retained by reference and
are unchanged (\artifact{dn-plane-principals}, \S 1--\S 2). Chapter 1's citations
to \artifact{dn-ouroboros-principal} should be read in this light: the reasoning
stands, the conclusion widened.
\end{devolution}

\subsection*{Four principals}

\phuman{} is the human: hand edits, blessings, the decisions. \pwork{} is the
workflow plane --- the orchestrator and its delegated builders and scribes, the
repository working tree, and the owner of the build-report lane. \pcore{} is the
core plane: the reasoning daemon and the vault, with read/write on the corpus and
no off-host network. \pedge{} is the edge plane: network-facing only, never the
vault (\artifact{dn-plane-principals}, \S 3.1).

The fourth is worth dwelling on, because it was created before it was needed. Edge
is effectively empty today --- the maximum reachable effector tier is
\textsc{none}, and the only edge surface is a Tailscale-local ambassador
(\artifact{finding-0011}) --- and the note creates the principal anyway, as
forward provisioning. The stated reason is the good one: ``it guards little now;
the model is complete the moment edge grows, and \emph{no future edge component can
be born inside core for convenience}'' (\artifact{dn-plane-principals}, \S 3.1).
This is a boundary drawn to remove a future temptation rather than a present
capability, which is an unusual and rather disciplined thing to spend a design on.

\subsection*{Non-negotiable \#2 becomes a kernel fact}

The constitutional line is that private data and network access never live in the
same component (\coderef{CONSTITUTION.md}{\gitref}, \S II.1). Before the plane
split, that was enforced at the \emph{source} level by the import firewall of
\autoref{sec:firewall} --- which is, as \artifact{dn-plane-principals} \S 3.4 puts
it precisely, review-time enforcement of a runtime property. The four-user split
makes it a kernel fact, enforced from both directions at once:

\begin{itemize}[leftmargin=*]
  \item \textbf{Edge cannot reach the vault.} $\pedge{} \neq \pcore{}$, and the
    vault is mode \texttt{0700} owned by \pcore{}. The kernel's discretionary
    access check denies the \texttt{open}. No code path, no import, and no bug in
    edge can cross it short of a kernel exploit or root.
  \item \textbf{Core cannot reach the network.} A \texttt{pf} anchor keyed on the
    uid drops core's off-host traffic, with a loopback carve-out placed
    \emph{before} it, because core legitimately speaks localhost to the resident
    model servers. ``Sealed'' means zero \emph{off-host} egress, not zero loopback.
\end{itemize}

The whole of the second bullet is two lines of configuration, and they are worth
reading in full because the ordering is the design:

\begin{quote}\small\ttfamily
source: ops/network/ouroboros-egress.pf.conf@\gitref
\end{quote}
{\small\begin{verbatim}
pass  out quick on lo0
block drop out quick proto { tcp udp } from any to any user ouroboros
\end{verbatim}}

\noindent \texttt{pf} is first-match-wins for \texttt{quick} rules, so the carve-out
must precede the block or the daemon's own model calls are dropped and core wedges.
The same file records that \texttt{pf}'s \texttt{user} clause matches the socket's
effective uid in \emph{both} directions, so this one block covers egress and any
listening socket core might open off-loopback. Two lines; one non-negotiable; a
kernel that does not negotiate.

\begin{figure}[t]
  \centering
  % figures/plane-principals.tex — one OS user per plane, and the two enforcement
% directions that compose into non-negotiable #2.
% Source: \artifact{dn-plane-principals} §3.1 (the ownership/mode matrix) and §3.4
% (the two directions: edge cannot reach the vault via DAC on a 0700 vault; core
% cannot reach the network via a pf anchor keyed on uid, with the lo0 carve-out).
% The pf rules are \coderef{ops/network/ouroboros-egress.pf.conf}{\gitref}:43--44.
% Principal names from notation.tex.
% \input by chapters/02-architecture.tex inside a figure environment.
\begin{tikzpicture}[
    node distance=9mm and 13mm,
    plane/.style={draw=inkblue, thick, rounded corners=3pt, fill=panel,
      align=center, inner sep=4pt, text width=34mm, minimum height=17mm,
      font=\small\sffamily},
    asset/.style={draw=soft, thick, fill=white, align=center, inner sep=3pt,
      text width=26mm, font=\scriptsize\sffamily},
    deny/.style={-{Bar[width=2mm]}, very thick, gate},
    allow/.style={-{Stealth[length=2mm]}, thick, inkblue},
    lbl/.style={font=\scriptsize\itshape, align=center, text=gate},
  ]
  \node[plane] (human) {\phuman{} --- the human\\[2pt]
    \scriptsize hand edits · blessings · the one who decides};
  \node[plane, below=of human] (work) {\pwork{} --- workflow\\[2pt]
    \scriptsize orchestrator + delegated builders; owns \texttt{exhaust/reports/}};
  \node[plane, right=32mm of work] (core) {\pcore{} --- core\\[2pt]
    \scriptsize the reasoning daemon + the vault};
  \node[plane, above=of core] (edge) {\pedge{} --- edge\\[2pt]
    \scriptsize network-facing only (forward-provisioned)};

  \node[asset, below=7mm of core] (vault) {the vault (corpus)\\\texttt{0700 ouroboros}};
  \node[asset, above=7mm of edge] (net) {the network};

  \draw[allow] (human) -- node[font=\scriptsize\sffamily, right=1mm, text=soft]
    {\texttt{sudo -u} (descending)} (work);
  \draw[allow] (core) -- (vault);
  \draw[allow] (edge) -- (net);

  % The two denials that compose.
  \draw[deny] (edge.south east) .. controls +(18mm,-10mm) and +(18mm,10mm) ..
    (vault.north east)
    node[lbl, right=1mm, midway, text width=30mm]
    {kernel DAC:\\ \pedge{} $\neq$ \pcore{}};
  \draw[deny] (core.north west) .. controls +(-9mm,9mm) and +(-9mm,-9mm) ..
    ([xshift=-3mm]net.west)
    node[lbl, left=1mm, midway, text width=27mm]
    {pf anchor on uid\\ (\texttt{lo0} carved out first)};

  \node[lbl, below=4mm of vault, text width=95mm, text=inkblue]
    {The component holding private data has no network; the component holding the
     network has no private data --- and they are \emph{different uids}, so no single
     process can hold both.};
\end{tikzpicture}

  \caption[The four plane principals]{One OS user per plane, and the two denials
    that compose. Source: \artifact{dn-plane-principals}, \S 3.1 (the ownership
    matrix) and \S 3.4 (the two directions); the rules are
    \coderef{ops/network/ouroboros-egress.pf.conf}{\gitref}.}
  \label{fig:planes}
\end{figure}

The composition is the payoff, and the note states it in one sentence: ``the
component holding private data has no network; the component holding the network
has no private data; and they are different uids, so no single process can ever
hold both'' (\artifact{dn-plane-principals}, \S 3.4). Note also what does
\emph{not} happen next: the import firewall is not retired. It stays as the cheap
pre-runtime tripwire --- ``it catches the violation at review, the kernel catches
it at execution.''

\subsection*{What the agents lose, and what the split does not do}

The clearest way to feel a capability boundary is to state it as a subtraction.
Running as \pwork{}, an agent --- and anything it spawns, and any supply-chain
compromise inside its toolchain --- cannot read the human's documents, SSH keys,
keychain or browser profiles; cannot open the vault; and cannot write core's
runtime substrate. Its writable world is the repository, its own home, and the
report lane. ``That \emph{is} the workflow plane, now enumerable by \texttt{ls
-l}'' (\artifact{dn-plane-principals}, \S 3.2).

And the residuals, stated by the note itself rather than discovered by a reader.
The split does \emph{not} enforce the blessing gates: the repository is
group-writable by design, so an errant agent process can still touch a ratified
note at the filesystem level --- the guards of \autoref{sec:workflow} remain the
enforcement there. Workflow keeps open network egress, because building needs it;
a workflow compromise can still exfiltrate the \emph{repository}, just no longer
the vault or the human's files. And \pwork{} must be able to read the human's
signing key in order to sign commits as the human --- a modest, knowingly accepted
leak, recorded with its rationale and with the alternative (a separate verified bot
identity) parked rather than discarded (\artifact{dn-plane-principals}, \S 3.2,
\S 5).

This is designed, built and instrumented rather than aspirational --- and the
instrument is the interesting part. The migration itself is owner-run and
privileged (creating users, loading a packet-filter anchor), so what ships in the
repository is a runbook, a set of per-plane ratchets, and a \emph{read-only}
verifier that mutates nothing and asserts the end state across all four principals
(\coderef{scripts/verify\_planes.py}{\gitref}). Crucially it is runnable
\emph{mid}-migration: a lane not yet migrated reports \textsc{pending}, never a
false green. A boundary whose instrument can say ``not yet'' is a boundary you can
trust the report of.

What the edge plane will eventually \emph{contain} is a separate question, and its
design note is not yet ratified --- so this manual names
\fwdthesis{the edge model}{ch:architecture} and asserts nothing about it beyond the
principal that already exists.

\section{Rings: a boundary of meaning, computed rather than curated}
\label{sec:rings}

Every boundary so far exists because something is dangerous. The next one exists
because something is \emph{confusing}, and it is the most instructive of the five,
because it shows the house style applied where there is no threat at all.

\subsection*{The motivation}

\texttt{core/} was heterogeneous: a pure mathematical and algebraic kernel, a
boundary/invariant layer, and administrative machinery, all in one package
(\artifact{dn-inner-outer-core}, \S 1). The owner's decision was to split it into
two rings --- but the interesting part is the constraint placed on \emph{how}: the
split had to be \textbf{mechanical}. Membership is computed, never curated, so that
no module's assignment is ever a judgment call.

The reason to insist on this is not tidiness. A curated boundary decays: each
individual exception is defensible, the accumulated set is not, and nobody can say
at any moment what the boundary now means. A computed boundary cannot decay,
because there is nothing to exercise judgment on. If you want to move it, you must
amend the \emph{predicate} --- a single ratified design decision, applied to
everything at once.

\subsection*{The predicate, stated exactly}

Let $\coremods{}$ be the modules under \texttt{core/} and let the admissible base be
\[
  \basemods{} \;=\; \bigl(\text{stdlib} \setminus \netmods{} \setminus
  \{\texttt{sqlite3}\}\bigr) \;\cup\; \{\texttt{numpy}, \texttt{scipy}\},
\]
where $\netmods{}$ is the same audited network-module set the firewall of
\autoref{sec:firewall} uses --- one home for ``network-capable,'' reused rather
than re-declared. Call a subset $S$ of $\coremods{}$ \emph{import-closed} when
every $m \in S$ satisfies $\imports(m) \subseteq S \cup \basemods{}$. Then

\begin{invariant}[The inner ring]
\label{inv:ring}
$\innerring{}$ is the maximal import-closed subset of $\coremods{}$ over
$\basemods{}$ (\artifact{dn-inner-outer-core}, \S 2.1, \S 2.3).
\end{invariant}

A maximum exists, and the reason is worth one line because it is what makes
``maximal'' well defined rather than a hopeful adjective: a union of import-closed
sets is import-closed (if $m \in S_1$ then $\imports(m) \subseteq S_1 \cup
\basemods{} \subseteq \bigcup_j S_j \cup \basemods{}$), so
$\innerring{} = \bigcup \{ S \subseteq \coremods{} : S \text{ import-closed}\}$ is
itself import-closed and contains every other. Equivalently, $\innerring{}$ is the
greatest fixed point of the monotone operator
$F(S) = \{ m \in \coremods{} : \imports(m) \subseteq S \cup \basemods{}\}$ on the
powerset lattice of $\coremods{}$ --- the classical lattice-theoretic
fixed-point construction.\footnote{A.~Tarski, ``A lattice-theoretical fixpoint
theorem and its applications,'' \emph{Pacific Journal of Mathematics} 5(2):285--309,
1955. \textsc{doi}:\,\texttt{10.2140/pjm.1955.5.285}. Verified 2026-07-25.} Two
immediate corollaries: an import-closed set cannot reach its complement, so
\emph{inner never imports outer} --- the direction law is a theorem, not a further
rule; and the fixed point is computable by iterating $F$ downward from
$\coremods{}$ until it stabilises, which is precisely what the ratchet test does.

\subsection*{Why exactly two subtractions --- one forced, one ruled}

The base could have been simply ``stdlib plus numpy and scipy.'' It is not, and the
record is unusually clear about the difference between the two corrections it makes
(\artifact{dn-inner-outer-core}, \S 2.1).

The first subtraction is \emph{computed-forced}. Applied literally, the naive base
lands the loopback model client (which imports \texttt{urllib}) and the egress guard
(which imports \texttt{socket}) \emph{inside} the inner ring, because
\texttt{urllib} and \texttt{socket} are stdlib. ``An inner core
containing the repo's only two network-capable core files would poison the term on
day one.'' Subtracting the audited network set makes \emph{inner $\subseteq$
network-incapable} a theorem with an empty allowlist --- and it does so by reusing
the firewall's constant rather than minting a second opinion about what counts as
network-capable.

The second subtraction is \emph{owner-ruled}, and is recorded as taste with its
reasoning attached rather than disguised as necessity. With only the first
subtraction the ring computes to include the SQLite store layer, because
\texttt{sqlite3} is stdlib and the store modules are import-austere. The ruling:
``the inner ring should \textbf{be} the founding language --- the algebra, the
mathematics, the sacred boundaries --- not the austere plumbing. Persistence
machinery, however dependency-clean, is machinery.'' The rejected wider alternative
is kept on the record with its own re-entry condition, not deleted.

\subsection*{The mechanism: a declaration forced to equal a computation}

Here is the move that makes the whole thing work, and it generalises far beyond
this chapter. The ring map is a plain declaration in code: a frozen set of module
names, living in
\coderef{core/kernel/rings.py}{\gitref}.
The ratchet test recomputes the fixed point at test time and asserts
\emph{equality in both directions}
(\artifact{dn-inner-outer-core}, \S 2.4):

\begin{quote}\small\ttfamily
pseudo-code --- the shape of tests/unit/test\_inner\_ring.py
\end{quote}
{\small\begin{verbatim}
computed <- greatest S subset of core/** with imports(S) subset of S + BASE
assert computed == DECLARED        # both directions, not containment
  # a module that acquires an inadmissible import LEAVES computed
  #   -> red until the map records the demotion, explicitly
  # a module that becomes pure ENTERS computed
  #   -> red until the map claims it, explicitly
\end{verbatim}}

\noindent Neither direction is decoration. Containment in one direction alone would
let the ring silently shrink; containment in the other would let it silently grow.
Equality means every change of membership is a \emph{reviewable one-file diff} that
must be named in a build plan's write scope --- so the guard of
\autoref{sec:workflow} makes every promotion and demotion plan-visible with no new
machinery at all. The design note is explicit that this is the point: the map ``is
a declaration forced to match a computation.'' And the failure mode is pre-assigned
--- if the ratchet goes red, the \emph{map} is stale, never the test; recompute at
\textsc{head} and edit the map, never hand-edit toward green.

Around it sit four no-laundering bindings, which exist because the interesting
attack on a ratchet is not to violate it but to redefine what it counts: no module
ever leaves \texttt{core/} under this programme; every physical move asserts the
outer violation count is unchanged on both sides; ring assignment cannot redefine a
violation (a violating module is outer \emph{because} it violates --- there is no
assignment step a violation could hide in); and the map may shrink only through a
named plan line-item (\artifact{dn-inner-outer-core}, \S 2.4).

\subsection*{Where it stands, and what it is not}

At the edition's ref the ring is 43 modules, and it physically lives in one
subtree, \texttt{core/kernel/**}
(\coderef{core/kernel/rings.py}{\gitref}).
The outer ring is the complement --- \texttt{core/**} minus that subtree --- and the
sibling packages
--- \texttt{eval}, \texttt{ops}, \texttt{agents}, \texttt{scheduler},
\texttt{edge}, \texttt{config} --- are not core at all
(\artifact{dn-inner-outer-core}, \S 2.5, \S 2.7). \autoref{fig:core-rings} draws
the three regions.

\begin{figure}[t]
  \centering
  % figures/core-rings.tex — the three-zone code picture, with the two boundaries
% drawn as what they actually are: the SECURITY perimeter (core-wide) and the
% MEANING boundary (the kernel ring inside it).
% Source: \artifact{dn-inner-outer-core} §2.5 ("inner ring subset core; outer ring =
% core minus inner; the siblings are not core at all" — and the honest correction:
% "the security perimeter is still core-wide ... the ring is a meaning boundary, not
% the egress boundary"), §2.7 (the physical end-state `core/kernel/**`).
% The perimeter is \coderef{ops/import\_lint.py}{\gitref} (all of core/**).
% \input by chapters/02-architecture.tex inside a figure environment.
\begin{tikzpicture}[
    lbl/.style={font=\scriptsize\sffamily, align=center},
    note/.style={font=\scriptsize\itshape, align=left, text=soft},
  ]
  % Outermost: the repo's sibling packages — not core at all.
  \node[draw=soft, thick, dashed, rounded corners=4pt, fill=white,
        minimum width=94mm, minimum height=54mm] (repo) {};
  \node[lbl, text=soft, anchor=north west] at ([shift={(2mm,-2mm)}]repo.north west)
    {\texttt{eval/} · \texttt{ops/} · \texttt{agents/} · \texttt{scheduler/} ·
     \texttt{edge/} · \texttt{config/}\\[1pt]
     \itshape the siblings --- machinery, not core};

  % The security perimeter: all of core/**.
  \node[draw=gate, very thick, rounded corners=4pt, fill=panel,
        minimum width=66mm, minimum height=36mm] (core) at (0,-4mm) {};
  \node[lbl, text=gate, anchor=north] at ([yshift=-1.5mm]core.north)
    {\textbf{the security perimeter} --- all of \texttt{core/**}\\[1pt]
     \itshape no network import, no \texttt{edge}/\texttt{cloud} import};

  % The meaning boundary: the kernel ring.
  \node[draw=inkblue, very thick, rounded corners=3pt, fill=white,
        minimum width=44mm, minimum height=17mm] (kernel) at (0,-9mm) {};
  \node[lbl, text=inkblue] at (kernel.center)
    {\textbf{the inner ring} --- \texttt{core/kernel/**}\\[1pt]
     \scriptsize the algebra · the mathematics ·\\[-1pt]
     \scriptsize the boundary vocabulary};

  \node[lbl, text=gate, anchor=south] at ([yshift=1.5mm]core.south)
    {the \textbf{outer} ring = \texttt{core/**} $\setminus$ \texttt{core/kernel/**}
     --- the machinery};

  % The direction law.
  \node[note, anchor=north west, text width=88mm]
    at ([shift={(3mm,-4mm)}]repo.south west)
    {Membership is \emph{computed}, never curated: the inner ring is the maximal
     import-closed subset of \texttt{core/**} over the admissible base. Outer imports
     inner, never the reverse --- a corollary of import-closedness, and given its own
     named test anyway.};
\end{tikzpicture}

  \caption[The three code regions]{Inner ring $\subset$ core; outer ring = core
    $\setminus$ inner; the siblings are not core at all. The security perimeter is
    drawn \emph{core-wide}, deliberately larger than the ring. Source:
    \artifact{dn-inner-outer-core}, \S 2.5 and \S 2.7.}
  \label{fig:core-rings}
\end{figure}

Two honest edges, both stated by the source rather than found by a critic. First,
\textbf{the ring is not the security perimeter}. Shrinking it surrendered nothing,
because the import firewall bans network imports across \emph{all} of
\texttt{core/**} regardless of ring: ``the ring is a meaning boundary, not the
egress boundary.'' Second, \textbf{inner does not mean pure}. Two members prove it
by inspection --- the content-addressed raw store performs disk I/O, and the config
loader reads the environment --- so ``nobody may build on \emph{inner $\Rightarrow$
pure function}.'' A call-grain purity predicate is parked, with a named candidate
consumer and a falsifier, rather than quietly assumed.

What the ring \emph{is}, the note says best, and it is the sentence to carry out of
this section: the inner core ``is not an agent and not a service --- it is the
vocabulary agents are written in; nobody grants you arithmetic''
(\artifact{dn-inner-outer-core}, \S 2.5).

\section{Types: the boundary at authorship time}
\label{sec:types}

The boundaries so far constrain \emph{reach}. This one constrains
\emph{expressibility} --- it moves the wall to the moment the code is written, so
that a whole class of wrong operation cannot be said.

\subsection*{What a green run actually proves}

The temptation with a type checker is to treat a green run as evidence the system
is correct. The record refuses that reading immediately, on the strongest possible
grounds. ``Under Curry--Howard a well-typed program \emph{is} a proof --- of
exactly the proposition its type expresses, and nothing
more''\footnote{The correspondence between proofs and programs, and a readable
modern account of it: P.~Wadler, ``Propositions as types,'' \emph{Communications of
the ACM} 58(12):75--84, 2015. \textsc{doi}:\,\texttt{10.1145/2699407}. Verified
2026-07-25.} (\artifact{dn-type-system-as-core-audit}, \S 2.1). The note makes the
point with a two-line counterexample:

\begin{quote}\small\ttfamily
pseudo-code --- the illustration of dn-type-system-as-core-audit \S 2.1
\end{quote}
{\small\begin{verbatim}
function add(a: int, b: int) -> int:
    return a - b          # type-checks.
    # The theorem proven is "ints map to ints" -- and that theorem is true.
\end{verbatim}}

\noindent So correctness splits into two obligations, and the split is the
section's central idea: \textbf{code $\models$ type-spec} is owned by the checker
--- mechanical and continuous; \textbf{type-spec $\models$ intent} is never
provable statically, and is owned by tests and by the owner. Knowing which of the
two a green run discharges is the difference between a useful instrument and a
false comfort.

There is a second boundary, from gradual typing, and it is sharper than it looks.
The guarantee holds only inside the \emph{checked region}. Unannotated functions
are silently exempt; \texttt{Any} is infectious, so a single \texttt{Any} at a
third-party boundary launders everything downstream out of the guarantee; casts
and suppressions are holes by construction. The consequence, stated flatly: ``a
green \emph{non-strict} run over partially annotated code is close to
meaningless.'' Strict mode earns its place not by proving more but by converting
silent exemptions into loud ones.

\subsection*{The design move: raise the propositional content}

If a type's value is the proposition it expresses, then the engineering lever is to
make the propositions say more. The note's formulation is the one to remember:
\texttt{int -> int} \emph{decorates}; a signature like

\begin{quote}\small\ttfamily
pseudo-code --- the static shadow of label monotonicity (\S 2.4)
\end{quote}
{\small\begin{verbatim}
promote(x: Derived[T], cap: OwnerVerdict) -> Authored[T]
\end{verbatim}}

\noindent \emph{constrains}. Under this signature, promoting derived material to
authored status without an owner verdict is not a policy violation to be detected
--- it is a type error at every call site, on every builder run, at zero runtime
cost. The same encoding gives capability non-amplification: a privileged operation
whose signature demands an unforgeable capability object simply cannot be called
from a path that never received one, and the checker verifies the plumbing end to
end (\artifact{dn-type-system-as-core-audit}, \S 2.4).

This is the standing razor of the whole project applied to annotations themselves:
\emph{formalism must constrain, not decorate}. An annotation earns its place by
encoding an invariant whose violation the checker would catch.

And --- characteristically --- the note immediately states how much weaker the
shadow is than the invariant it shadows. It sees no values and no runtime paths,
and a deliberate cast defeats it. What it removes is the \emph{accidental}
violation class, at authorship time. That is exactly the aligned move: ``the
correct design removes the error class rather than detecting instances of it.''

\subsection*{The audit framing, and a boundary that grows by itself}

The first strict pass over \texttt{core/} was reframed as an \emph{audit}, whose
error inventory is a measurement rather than a chore
(\artifact{dn-type-system-as-core-audit}, \S 2.3; warrant \artifact{finding-0026}).
Each error triages into exactly one of three classes: \textbf{T1}, a latent defect
corresponding to reachable incorrect behaviour, each filed as its own finding;
\textbf{T2}, a \emph{representability} finding --- behaviourally fine, but the
types reveal untyped shape crossing a module boundary; \textbf{T3}, checker
friction, resolved by a suppression that must carry an error code and a warrant
comment so that bare suppressions are grep-detectable violations.

The T2 class is the one worth internalising, because it is the type system doing
boundary work rather than correctness work: \emph{a shape the type system cannot
see is a shape the code-plane audit cannot defend}. T2 density measures how much of
the system's real structure lives outside the checked region.

The scope rule is the owner's, and it has the same computed-not-curated character
as the ring. Because a checker verifies call sites in the \emph{caller's} module,
``a caller must respect the callee's types'' is enforceable only if every caller of
core is itself analysed. Hence two tiers: core plus the foundation file set under
full strict, and an interface tier consisting of \emph{every module that imports
anything from} \texttt{core/}. Membership of the second is ``not a judgment call
but a mechanical invariant --- \texttt{imports core} $\Rightarrow$ present in
Tier-2 config --- decidable from the AST import graph and enforceable on the gate
path,'' and the ratchet is one-way: no module ever moves to a weaker tier, and
Tier-2 membership ``grows automatically with the import graph rather than by
intention'' (\artifact{dn-type-system-as-core-audit}, \S 2.5). A boundary that
enlarges itself as the code changes needs no maintainer.

Where an external dependency ships no honest types, the \texttt{Any} is quarantined
in a thin typed wrapper --- one file per dependency, rather than smeared through
core. That layer is built: \coderef{core/typedshims/lancedb.py}{\gitref} is, in its
own words, ``the ONE place core touches the raw package,'' pinning returned values
to the minimal protocol surface the store actually uses, with the single
suppression carrying its warrant inline. Its final instruction is a small lesson in
boundary hygiene: do not widen these protocols speculatively --- each addition
should arrive with the call that needs it.

Finally, the falsifier, because this chapter would be dishonest without it. The
note names the observation that would retire its own central claim: if the first
strict pass yields zero T1 and zero T2 findings and the steady-state cost is
dominated by T3 friction, then the \emph{audit} framing has failed and is recorded
as no-signal --- the checker may survive as a regression floor, but the stronger
claim is withdrawn (\artifact{dn-type-system-as-core-audit}, \S 2.2).

\section{Scope: the boundary at query time}
\label{sec:scope}

Every boundary so far is a fact about the system --- true of the code, the
processes, the users. This one is a fact about a \emph{request}: what a given
caller, right now, is permitted to ask.

It is also, unlike the others, a piece of \emph{vocabulary} rather than a gate. The
built module says so in its own header: it ``wires NO enforcement into any read
path --- it is vocabulary, not a gate''
(\coderef{core/kernel/scope.py}{\gitref}). Read this section as the language in
which every other boundary in the chapter can be stated exactly, and in which future
grants will be checked.

\subsection*{The object}

A scope is a four-tuple $\scopeobj{} = (\stratums{}, \edgeclass{}, \timescope{},
\authority{})$ (\artifact{dn-capability-scope}, \S 2.1). Informally: \emph{which
material}, \emph{which relations over it}, \emph{as of when}, and \emph{with what
authority}. Each component is a lattice, and the interesting content is in how each
is made well-posed.

\paragraph{$\stratums{}$ --- which material.} A downward-closed subset of a
refinement forest $\reffor{}$ whose base strata are the system's kinds of material
--- mirror, curated, observed, ops, reference, interpreted, world
(\artifact{dn-capability-scope}, \S 2.1), joined since by a \emph{dialogue}
stratum --- with refinement predicates below them as first-class \emph{elements}
rather than annotations. The refinements are where the boundary work happens: the
built forest carries \texttt{exhaust} as a refinement of dialogue that is
\emph{default-excluded} and grantable only when named
(\coderef{core/kernel/scope.py}{\gitref};
\artifact{dn-agentic-loop}, \S 2.4b) --- \autoref{sec:workflow}'s one-way exhaust
invariant, re-expressed as a lattice element. The detail that matters most here is
the top:
\[
  \top_{\stratums{}} \;=\; \reffor{} \setminus \denyset{},
\]
where $\denyset{}$ is the foundation denylist. \emph{Even the fullest grant
structurally excludes} \texttt{CONSTITUTION.md} \emph{and the frozen golden set}.
The fixed point of \autoref{sec:fixed-point} is not protected by a check that runs
before each write; it is subtracted from the space of grantable scopes, so no
grant that includes it can be constructed.

\paragraph{$\edgeclass{}$ --- which relations.} A subset of the edge fibers.
The algebra was ratified with two --- $F$, citation-support, and $D$, supersession
lineage --- and the agent taxonomy argued for a third, $C$ (causal-witnessed
\emph{origin}), as an additive lattice extension. The argument for why $C$ is not
$D$ is worth the detour, because it is a boundary drawn between two things that look
alike: a supersession edge is \emph{homogeneous}, two versions of one artifact's
lineage, whereas ``this conversation produced this change'' is
\emph{heterogeneous} --- origin, not lineage. ``An utterance is not a prior version
of a file'' (\artifact{dn-agent-taxonomy}, \S 2.5). The extension has since landed:
the built algebra types $\edgeclass{} \subseteq \{F, D, C\}$
(\coderef{core/kernel/scope.py}{\gitref}).

\paragraph{$\timescope{}$ --- as of when.} A pair (clock $\clockk{}$, window),
where a clock is a monotone coarsening of the append-only ledger's causal order.
This component carries the chapter's best example of a boundary drawn around the
system's own ignorance. Clocks form a hierarchy, and the finest --- a global event
clock over all strata --- \emph{does not yet exist}. So the meet on
$\timescope{}$ is honestly partial: same clock, intersect the windows; different
clocks, pull both windows back to a materialised common refinement; and if no such
clock has been materialised, the constructor \textbf{raises an error} ---
``\emph{never a silent guess}'' (\artifact{dn-capability-scope}, \S 2.2). A
lattice operation that refuses to be total is an unusual thing to design on
purpose. It is the alternative to a system that quietly compares two timelines that
were never comparable.

\paragraph{$\authority{}$ --- with what power.} A product of three chains,
\[
  \authority{} \;=\; P \times W_{\stratums{}} \times W_{\text{world}},
\]
being the advisory ladder $P \in \{\textsf{read} < \textsf{read+propose}\}$ ---
which is non-negotiable \#3, ``the model advises; code acts,'' as a lattice
coordinate --- the projection-write rung $W_{\stratums{}}$, and the effector
blast-radius chain $W_{\text{world}}$. The product exists to repair a real
conflation: a single ``write'' rung had merged \emph{sensor projection} (writing
into the store) with \emph{effector mutation} (acting on the world), ``two
capabilities the architecture already separates structurally --- a sensor has no
world reach, the executor no store reach'' (\artifact{dn-capability-scope},
\S 2.1). Track G's standing fact --- that the maximum reachable effector tier is
\textsc{none} --- becomes a single lattice statement:
$\top_{\text{deployed}}.W_{\text{world}} = \textsc{none}$.

That three-way split of authority is the \emph{ratified partial} of a larger
argument about what crosses the core's boundary. The fuller taxonomy ---
\fwdthesis{the sacred boundary}{ch:architecture} --- is not yet ratified, so this
manual names it and withholds its content; \artifact{finding-0180} records the gap
and proposes the note for ratification.

\subsection*{Composition, and delegation as a law}

Meet is componentwise, and meet is \emph{safe composition}: a delegated agent
receives the meet of its parent's scope with its template's. That single equation is
non-negotiable \#6 --- you cannot delegate more than you hold --- expressed as
arithmetic rather than as an instruction
(\artifact{dn-capability-scope}, \S 2.2, \S 3):
\[
  \scopeobj{}_{\text{minted}} \;=\; \scopeobj{}_{\text{parent}} \meetop{}
  \scopeobj{}_{\text{template}}
  \qquad\text{hence}\qquad
  \scopeobj{}_{\text{minted}} \scopeleq{} \scopeobj{}_{\text{parent}}.
\]
Monotone delegation is then a theorem about a meet-semilattice rather than a rule
an agent is trusted to obey. \autoref{sec:structural} already showed the
code-side floor of this at work: the frozen one-element maximum that construction
refuses to exceed
(\coderef{core/factory/roles.py}{\gitref}).
The algebra is what that floor is a special case of.

Join is a widening, grantable only by an authority that already holds the join.
And enforcement tier is carried as an \emph{annotation}, never as a lattice
element, with composition taking the minimum along the chain --- the law of
\autoref{sec:tiers}, which is where it comes from.

\subsection*{Firewalls as ideals: ``dissolve, do not guard'' as lattice theory}

Here is the moment where \autoref{sec:sacred}'s principle stops being a slogan.

Each firewall --- the mirror-payload firewall for non-exempt clients, the
foundation denylist always --- is an \emph{order ideal} $\iota$ in the scope
lattice. A grant is admissible for a client class $c$ exactly when
\[
  \scopeobj{} \meetop{} \iota \;=\; \bot
  \qquad \text{for every } \iota \text{ applicable to } c ,
\]
and the grantable lattice is therefore the ideal quotient. The design note's own
gloss is the sentence to hold onto: \emph{firewalls are subtracted from the space
of scopes, not re-checked per query} (\artifact{dn-capability-scope}, \S 2.2).

Read that against Chapter 1's principle. ``Dissolve, do not guard'' says: remove
the need for the dangerous capability rather than permitting it under controls.
Here the dangerous capability is a scope that reaches forbidden material, and it is
not refused at query time --- it is \emph{not in the lattice}. There is no query to
refuse because there is no scope to construct. The principle and the algebra are
the same statement in two languages, which is the strongest evidence available that
the principle was a real one.

The same move governs the query language above it. A query is a sentence
$(\text{verb}, \scopeobj{})$, admissibility is
$\text{req}(\text{verb}) \scopeleq{} \scopeobj{}_{\text{granted}}$ checked at
construction, and ill-scoped sentences are \emph{unrepresentable}. And mode is a
corollary of scope rather than a flag on a call: whether a query is structural,
semantic or temporal follows from which fibers and which window the scope carries
(\artifact{dn-capability-scope}, \S 2.4).

\subsection*{Two consequences worth their own names}

\paragraph{The slice rule.} A scope spanning more than one stratum with a
point-in-time window must carry an explicit \emph{consistent cut}; a bare ``now'' is
well-typed only within a single stratum. This is the classical observation that an
instant is not globally defined across independent
timelines,\footnote{L.~Lamport, ``Time, clocks, and the ordering of events in a
distributed system,'' \emph{Communications of the ACM} 21(7):558--565, 1978.
\textsc{doi}:\,\texttt{10.1145/359545.359563}; K.~M.~Chandy and L.~Lamport,
``Distributed snapshots: determining global states of distributed systems,''
\emph{ACM Transactions on Computer Systems} 3(1):63--75, 1985.
\textsc{doi}:\,\texttt{10.1145/214451.214456}. Both verified 2026-07-25.} applied
to a corpus rather than to a network. For repo-backed strata the commit SHA
\emph{is} the consistent cut --- and the note observes, with evident pleasure, that
this retro-explains an existing piece of code: two independent views were already
routed through one shared commit resolver
(\coderef{core/reference\_view.py}{\gitref}, \texttt{\_resolve\_default\_commit}).
``The slice discipline was discovered empirically; the type now enforces it''
(\artifact{dn-capability-scope}, \S 2.2). A boundary can be found before it is
formalised; formalising it is what stops the next author from unfinding it.

\paragraph{$\Inv{}$ versus $\Rate{}(\clockk{})$.} Results are graded by
clock-dependence. An $\Inv{}$ result depends only on the window's event set ---
counts, sets, booleans. A $\Rate{}(\clockk{})$ result is a difference quotient
against a clock's index --- severings per commit, events per wall-second. The rule:
computing a rate against clock $\clockk{}$ \emph{requires} that the scope's own
clock be $\clockk{}$, and a rate value carries its clock in its type and is never a
bare number (\artifact{dn-capability-scope}, \S 2.3). The failure this forecloses
is precise and would otherwise be invisible: a drift \emph{rate} read off an
unacknowledged clock --- a number that looks like a measurement and is an artefact
of the denominator. Every instrument built so far audits as $\Inv{}$; the coherence
report, for instance, returns its count together with both anchors and deliberately
never divides.

\section{The workflow is a zone}
\label{sec:workflow}

Everything so far concerns the running system. The final boundary is the one that
governs the process that \emph{builds} it --- and treating that process as another
zone, with the same geometry and the same rubric, is one of the project's more
distinctive decisions.

It is also not merely an analogy. The instruments are literally shared. The
foundation denylist appears twice: as the order ideal $\denyset{}$ subtracted from
the top of the scope lattice (\autoref{sec:scope}), and as a list of paths refused
to every session by a hook. \autoref{ch:philosophy} established the artifact chain
and the plan-as-capability; this section gives the mechanism, and a reader who has
followed the chapter will recognise every move in it.

\subsection*{The plan is a capability, and the capability is checked twice}

A build plan's \texttt{write\_scope} is a list of globs, and it is enforced pre-hoc
on every edit (\artifact{dn-agent-workflow}, \S 6). The decision procedure is short
enough to state in full:

\begin{quote}\small\ttfamily
pseudo-code --- the shape of .claude/hooks/\_lib.py, cmd\_scope\_check
\end{quote}
{\small\begin{verbatim}
on every Edit/Write, given a target path p:
  if p matches FOUNDATION_DENYLIST:   DENY   # beneath any plan, any session
  if p is a design note, status in {ratified, superseded}:
                                      DENY   # the blessed record is immutable
  if no plan is active:               ALLOW  # orchestrator posture
  allowed <- plan.write_scope + [the plan, its journal, docs/findings/**]
  if p matches allowed: ALLOW else DENY
\end{verbatim}}

\noindent Three details in that sketch repay attention.

First, the denylist is \emph{location}-keyed and the design-note guard is
\emph{status}-keyed, and the asymmetry is deliberate. The Constitution and the
frozen golden set are fixed points --- they are what they are wherever they sit. A
design note's protection, by contrast, must follow its \emph{blessing}: a draft is
agent-writable working material, and the same file becomes agent-immutable the
moment the owner ratifies it. Guarding notes by location would either freeze
drafts or expose ratified text; guarding them by status tracks the thing that
actually changed (\artifact{dn-agent-workflow}, \S 6; \coderef{CLAUDE.md}{\gitref}).

Second, the always-appended tail --- the plan, its journal, and the findings
directory --- is the structural reason a constrained agent is never cornered. An
agent that discovers it needs to write outside its capability has exactly one
sanctioned move, and that move is always available: file a finding. The boundary is
tight precisely because there is a legitimate door.

Third, and most honest: this pre-hoc guard covers editing tools only. An agent can
write files through a shell, and pre-hoc shell pattern-denial is brittle and was
deliberately parked. So the enforcement is two-layer, and the record states the
posture in four words --- \emph{pre-hoc porous, post-hoc tight}
(\artifact{dn-agent-workflow}, \S 6). At session close, an audit diffs the working
tree and blocks the session from ending on any out-of-scope path, forcing revert
or finding first. The audit is untracked-inclusive, because a shell-written file is
precisely the untracked one the pre-hoc guard could not see.

The immutability check gets the same treatment, and the phrase for it is worth
borrowing: the pre-hoc arm reads status \emph{on disk}, the post-hoc arm reads
status \emph{at} \textsc{head} --- \emph{laundering-proof}, because an agent that
edits a ratified note and also rewrites its status line to \texttt{draft} still
fails the comparison against the committed truth.

\subsection*{The clause family}

The Stop audit is not one check but a family of six lettered clauses (one with a
sub-clause), added one at a time as each gap was found, each traceable to the
artifact that decided it:

\begin{itemize}[leftmargin=*,itemsep=1pt]
  \item \textbf{(a)} the journal has not been written since the last commit;
  \item \textbf{(b)} the uncommitted delta reaches outside the plan's write scope;
  \item \textbf{(b2)} a blessed design note differs from \textsc{head} (A8);
  \item \textbf{(c)} the working tree contains an uncommitted blessing transition;
  \item \textbf{(d)} cross-checkout state bleed between worktrees;
  \item \textbf{(e)} the session committed work but left the resume brief stale
    (\artifact{dn-session-handoff-gate});
  \item \textbf{(f)} a plan was sealed to \texttt{complete} without a
    \emph{follow-through} block in its journal
    (\artifact{dn-track-board-and-deskcheck-gate}, \S 2.5).
\end{itemize}

Clause (e) is a good miniature of the house method. The handoff loop is seal
$\to$ brief $\to$ kill $\to$ resume; three of the four steps were already
automated, and the brief was ``authored by habit, enforced by nothing''
(\artifact{dn-session-handoff-gate}, \S 1). Rather than add machinery, the note
found a signal already lying around --- the session's starting commit, written by
an existing hook --- and expressed the block in two lines:

\begin{quote}\small\ttfamily
restatement of \artifact{dn-session-handoff-gate} \S 2.2 --- state paths abbreviated
\end{quote}
{\small\begin{verbatim}
BLOCK  iff  HEAD != content(session-baseline)   # commits happened this session
       and  mtime(resume-brief) < time of last commit
\end{verbatim}}

\noindent and then --- the part that makes it trustworthy --- states its own
porosity: an \texttt{mtime} is launderable by a shell command; commit-less sessions
can still end briefless; a one-line brief passes because the gate checks
\emph{existence}, not quality. ``The gate targets forgetting, not adversarial
evasion'' (\artifact{dn-session-handoff-gate}, \S 2.6). Naming the threat model a
guard \emph{does not} cover is what keeps the guard from being mistaken for a wall.

The failure posture beneath all six is the same and is equally candid: hooks fail
open on script error and that platform behaviour is not configurable, so every
script runs under a trap that emits a conspicuous failure line and appends a marker
to the journal --- \emph{fail open, fail loud}
(\artifact{dn-agent-workflow}, \S 6). A guard that could fail silently would be
worse than no guard, because it would be believed.

\subsection*{Derivation instead of enforcement}

The strongest move in the workflow zone is the one that removes a guard entirely.

Project state --- which track is in which phase, which reviews are owed --- is not
a maintained file. It is \emph{derived}: recomputed from the artifact tree on every
run, with no persisted state, so it cannot drift
(\artifact{dn-track-board-and-deskcheck-gate}, \S 2.2). The consequence is stated
in one line that could serve as the chapter's epigraph: \emph{a board that cannot
be stale needs no hook to keep it honest.} The rejected alternative was a Stop
clause forcing agents to update the board at seal time --- described, correctly, as
``a hook where a compiler belongs.''

The same reasoning governs the third owner-only gate. A track is not closed by
setting a flag; closure is \emph{computed} --- a track is closed exactly when an
approved review record names it and its definition-of-done items are closed. ``No
artifact carries a \emph{closed} boolean \ldots{} nothing exists for an agent to
launder'' (\artifact{dn-track-board-and-deskcheck-gate}, \S 2.3). This is
\autoref{sec:scope}'s ideal-subtraction argument in a
different medium: rather than forbidding an action, arrange for the state it would
have produced to be unrepresentable.

\subsection*{The exhaust lane closes the loop}

One last boundary, and it returns the chapter to where \autoref{ch:philosophy}
began.

The system has an inbound lane --- the vault, owner-written and system-read --- and
the design gave it a mirror: an \emph{exhaust} lane, system-written and owner-read,
carrying artefacts the system emits for the owner to read away from the keyboard
(\artifact{dn-exhaust-lane}, \S 1). Two lanes, opposite directions, and one
invariant between them:

\begin{invariant}[Exhaust never re-enters]
\label{inv:exhaust}
No configured ingest source root may lie inside the exhaust lane, and no core
ingest path may read from it (\artifact{dn-exhaust-lane}, \S 2.2).
\end{invariant}

The reason is the ouroboros's reason. ``Reports describe the system; ingesting them
is recursive self-ingestion of meta-content.'' A system that indexes its own
descriptions of itself is not deepening its self-knowledge; it is amplifying its
own summaries until the summaries outweigh the material.

And the \emph{layout} was chosen to make the invariant structural rather than
observed. Exhaust is a sibling directory of the vault, not a subdirectory of it,
and the argument for that choice is the chapter's thesis in miniature: the ingest
scanner's source roots are config-driven, so a sibling directory is \emph{outside
every configured root by construction} --- the scanner has no path to it. The
unified layout ``would instead depend on an exclusion rule being honored forever.''
\emph{Structural isolation beats an honored exclude}
(\artifact{dn-exhaust-lane}, \S 2.1). The invariant is additionally held by a test
that loads the merged configuration and asserts no source root resolves inside
exhaust --- but the test guards a property the directory layout already makes
difficult, which is the correct order of operations.

So the serpent may consume its tail. It ingests its own reasoning, spawns its own
agents, and proposes changes to itself --- and there is one thing it may never
digest (the fixed point) and one lane it may never re-enter. That is the whole of
the architecture, if it must be compressed to a sentence.

\section{What the boundaries do not do}
\label{sec:residuals}

A chapter that only listed defences would be asking to be trusted, which is exactly
what \autoref{sec:falsifiers} says a claim may not do. Every residual below is
stated by the source that owns the mechanism, not discovered by an outside critic
--- and that provenance is itself the argument for believing the rest.

\begin{itemize}[leftmargin=*]
  \item \textbf{The import firewall} proves that exactly two audited files touch
    networking primitives, both loopback-only --- not that no core file imports
    \texttt{http}. The Ollama channel is real
    (\coderef{ops/import\_lint.py}{\gitref}). And, per
    \autoref{prop:closure-limit}, its \emph{global} no-egress-path claim is
    conditional on core self-containment, whose ratchet
    (\artifact{finding-0103}; \artifact{finding-0182}) is not yet at zero:\\
    \hspace*{1em}\coderef{tests/unit/test\_core\_self\_containment.py}{\gitref}.\\
    The runtime and kernel layers hold the property unconditionally meanwhile.
  \item \textbf{The runtime seal} is a Python-level guard; a native extension
    opening its own socket could bypass it. It fails closed and composes with the
    kernel layer, and claims no more (\coderef{core/sealing.py}{\gitref}).
  \item \textbf{The plane split} does not enforce the blessing gates --- the
    repository is group-writable by design --- and the workflow plane keeps open
    egress, so a workflow compromise can still exfiltrate the repository. \pwork{}
    can also read the human's signing key. All three are named and accepted, one
    with a parked alternative (\artifact{dn-plane-principals}, \S 3.2, \S 5).
  \item \textbf{The ring} is a meaning boundary, not the egress boundary; the
    security perimeter remains core-wide. And inner does not imply pure: two
    members perform disk and environment reads (\artifact{dn-inner-outer-core},
    \S 2.5).
  \item \textbf{The type shadow} is strictly weaker than the invariant it shadows:
    it sees no values and no runtime paths, and a deliberate cast defeats it. It
    removes the accidental violation class only
    (\artifact{dn-type-system-as-core-audit}, \S 2.4).
  \item \textbf{The scope algebra} is a \emph{typing layer}, not a gate --- the
    built module states that it wires no enforcement into any read path
    (\coderef{core/kernel/scope.py}{\gitref}). Its tier is per-inhabitant: two of
    the five built views are \tstruct{}, three are \tstatic{} and are labelled so
    (\artifact{dn-capability-scope}, \S 2.4). Its time component is honestly
    partial: cross-clock meets raise rather than guess.
  \item \textbf{The workflow guards} are porous pre-hoc by construction and tight
    only post-hoc; an \texttt{mtime} is launderable; the handoff gate checks that
    a brief exists, not that it is good. They target forgetting, not an adversary
    (\artifact{dn-agent-workflow}, \S 6; \artifact{dn-session-handoff-gate},
    \S 2.6).
\end{itemize}

\subsection*{The arguments still being developed}

Two pieces of this chapter's subject matter are, at this edition, designs in
progress, and the manual's rule is that a design in progress is named and never
asserted (\artifact{dn-agent-workflow}, \S 3).

The first is \fwdthesis{the sacred boundary}{ch:architecture} --- the full taxonomy
of the channels that cross the core's boundary, of which \autoref{sec:scope} gives
only the ratified fragment (the three-factor authority product). Chapter 1
originally promised that taxonomy here; the promise is now a forward reference, and
\artifact{finding-0180} records the gap and proposes the note for ratification.
The second is \fwdthesis{the edge model}{ch:architecture} --- what the edge plane
will contain once it has tenants; \autoref{sec:planes} gives only the principal
that already exists.

A third gap is of a different kind and is recorded here because it bears on how
this chapter should be read. The composition argument that would tie
\autoref{sec:firewall}, \autoref{sec:types} and \autoref{sec:scope} together ---
the claim that code-level, data-level and capability-level defences do not overlap
and do not substitute for one another --- lives in a research note that has not
been ratified. This chapter therefore presents the three as three mechanisms with a
shared shape, which is honest, rather than as a proven composition, which would not
be. \artifact{finding-0181} records it.

\bigskip

The boundaries are drawn. \autoref{ch:math} turns to the objects that live inside
them --- the coboundary framing and the instruments derived from it ---
\autoref{ch:intuition} assembles the whole into a working mental model, and
\autoref{ch:future} records what has been deliberately deferred, each parked
decision with the exact condition under which it returns.
